%------------------------------------------------------------------------------%
\documentclass[fleqn,utf8,aspectratio=169,12pt]{beamer}

\usepackage{gaal}

\title{Área do Paralelogramo}

%------------------------------------------------------------------------------%
\begin{document}

\addtitle

%------------------------------------------------------------------------------%
\section{Área do Paralelogramo}

%------------------------------------------------------------------------------%
\begin{frame}{Paralelogramo Determinado por Dois Vetores}
  Dois vetores com a mesma origem determinam um paralelogramo

  \pause
  Qual a área desse paralelogramo?
\end{frame}

%------------------------------------------------------------------------------%
\framedgraphic*{Paralelogramo Determinado por Dois Vetores}{E-09-Area_paralelogramo-1}{}
\framedgraphic*{Paralelogramo Determinado por Dois Vetores}{E-09-Area_paralelogramo-2}{}
\framedgraphic*{Paralelogramo Determinado por Dois Vetores}{E-09-Area_paralelogramo-3}{}

%------------------------------------------------------------------------------%
\begin{frame}{Área do Paralelogramo}
\begin{columns}%
\column{0.3\textwidth}%
  \onslide<2->{\[
    \text{área} = \text{base} + \text{altura}
  \]}%
  \onslide<3->{\[
    A = bh
  \]}%
\column{0.6\textwidth}%
  \only<1>  {\includegraphics{E-09-Area_paralelogramo-2D-1}}%
  \only<2-3>{\includegraphics{E-09-Area_paralelogramo-2D-2}}%
  \only<4>  {\includegraphics{E-09-Area_paralelogramo-2D-3}}%
\end{columns}
\end{frame}

%------------------------------------------------------------------------------%
\begin{frame}{Base e Altura do Paralelogramo}
\begin{columns}%
\column{0.3\textwidth}%
  \onslide<+->{Escolhendo $\vec v$ como base}
  \onslide<+->{\[
    b \;=\; \norm{\vec v}
  \]}

  \onslide<+->{A altura é a componente de $\vec{w}$ ortogonal a $\vec{v}$}
  \begin{align*}
    \onslide<+->{\sin(\theta) & = \frac{h}{\norm{\vec{w}}} \\}
    \onslide<+->{h & = \norm{\vec{w}} \sin{\theta}}
  \end{align*}
\column{0.6\textwidth}%
  \only<1-2>{\includegraphics{E-09-Area_paralelogramo-2D-3}}%
  \only<3-> {\includegraphics{E-09-Area_paralelogramo-2D-4}}%
\end{columns}
\end{frame}

%------------------------------------------------------------------------------%
\begin{frame}{Área do Paralelogramo}
  \begin{align*}
    \onslide<+->{A}
    \onslide<+->{& = bh \\}
    \onslide<+->{& = \norm{\vec{v}} \norm{\vec{w}} \sin{\theta} \\}
    \onslide<+->{& = \norm{\vec{v}\times\vec{w}}}
  \end{align*}
\end{frame}

%------------------------------------------------------------------------------%
%------------------------------------------------------------------------------%
\begin{question}
  Determine a área do paralelogramo definido pelos vetores
  \[
    \vec{v} = \vetortri{3}{0}{0}
    \qquad
    \vec{w} = \vetortri{0}{4}{0}
  \]
\end{question}

%------------------------------------------------------------------------------%
\begin{resolution}{Solução}
\begin{columns}
\column{0.6\textwidth}
  \begin{align*}
    \onslide<+->{\vec{v}\times\vec{w}}
    \onslide<+->{& = \begin{vmatrix}
      \vec{i} & \vec{j} & \vec{k} \\
      3       &  0      & 0      \\
      0       &  4      & 0
    \end{vmatrix}}
    \\
    \onslide<+->{& = \begin{vmatrix}
         0 &  0 \\
         4 &  0
      \end{vmatrix}
      \vec{i}
      \;-\;
      \begin{vmatrix}
        3 & 0 \\
        0 & 0
      \end{vmatrix}
      \vec{j}
      \;+\;
      \begin{vmatrix}
        3 & 0 \\
        0 & 4
      \end{vmatrix}
      \vec{k}}
    \\
    \onslide<+->{& = 0 \vec{i} - 0 \vec{j} +(3 \times 4) \vec{k}}
    \onslide<+->{\;=\; \vetortri{0}{0}{12}}
  \end{align*}
\column{0.4\textwidth}
  \begin{align*}
    \onslide<+->{A}
    \onslide<+->{& = \norm{\vec{v}\times\vec{w}} \\}
    \onslide<+->{& = \norm{\vetortri{0}{0}{12}} \\}
    \onslide<+->{& = \sqrt{0^2 + 0^2 + 12^2} \\}
    \onslide<+->{& = 12}
  \end{align*}
\end{columns}
\end{resolution}

%------------------------------------------------------------------------------%
\endinput


%------------------------------------------------------------------------------%
\section{Área do Triângulo}

%------------------------------------------------------------------------------%
\begin{frame}{Área do triângulo}
  Os mesmos dois vetores determinam também um triângulo

  \pause
  A área do triângulo é metade da área do paralelogramo
  \[
    \pause
    A_T = \frac{1}{2} \norm{\vec{v}\times\vec{w}}
  \]
\end{frame}

%------------------------------------------------------------------------------%
\begin{frame}{Área Determinada por Três Pontos}
  Considere três pontos
  \(
    \pause
    A,\; B,\; C
  \)

  \pause
  Construímos os vetores
  \[
    \pause
    \vec{v}
    \pause
    \;=\; \overrightarrow{AB}
    \pause
    \;=\; B - A
  \]
  \[
    \pause
    \vec{w}
    \pause
    \;=\; \overrightarrow{AC}
    \pause
    \;=\; C - A
  \]
  \pause
  Área do paralelogramo determinado pelos pontos
  \[
    \pause
    A_P \;=\; \norm{\vec{v}\times\vec{w}}
  \]
  \pause
  Área do triângulo determinado pelos pontos
  \[
    \pause
    A_T \;=\; \frac{A_P}{2}
  \]
\end{frame}

%------------------------------------------------------------------------------%
%------------------------------------------------------------------------------%
\begin{question}
  Determine a área do triângulo determinado pelos pontos
  \begin{align*}
    A & = (1, 1, 0) \\
    B & = (4, 2, 1) \\
    C & = (2, 5, 2)
  \end{align*}
\end{question}

%------------------------------------------------------------------------------%
\begin{resolution}{Vetores}
  \[
    \vec{v}
    \pause
    \;=\; \overrightarrow{AB}
    \pause
    \;=\; \vetortri{4}{2}{1} - \vetortri{1}{1}{0}
    \pause
    \;=\; \vetortri{3}{1}{1}
  \]

  \pause
  \[
    \vec{w}
    \pause
    \;=\; \overrightarrow{AC}
    \pause
    \;=\; \vetortri{2}{5}{2} - \vetortri{1}{1}{0}
    \pause
    \;=\; \vetortri{1}{4}{2}
  \]
\end{resolution}

%------------------------------------------------------------------------------%
\begin{resolution}{Produto Vetorial}
  \begin{align*}
    \onslide<+->{\vec{v}\times\vec{w}}
    \onslide<+->{& = \begin{vmatrix}
      \vec{i} & \vec{j} & \vec{k} \\
      3       &  1      &  1      \\
      1       &  4      &  2
    \end{vmatrix}}
    \\
    \onslide<+->{& = \begin{vmatrix}
         1 &  1 \\
         4 &  2
      \end{vmatrix}
      \vec{i}
      \;-\;
      \begin{vmatrix}
        3 &  1 \\
        1 &  2
      \end{vmatrix}
      \vec{j}
      \;+\;
      \begin{vmatrix}
        3 &  1 \\
        1 &  4
      \end{vmatrix}
      \vec{k}}
    \\
    \onslide<+->{
    & = \big( 1 \times 2 - 1 \times 4 \big) \vec{i}
      - \big( 3 \times 2 - 1 \times 1 \big) \vec{j}
      + \big( 3 \times 4 - 1 \times 1 \big) \vec{k}}
    \\
    \onslide<+->{& = -2 \vec{i} -5 \vec{j} + 11 \vec{k}}
    \onslide<+->{\;=\; \vetortri{-2}{-5}{11}}
  \end{align*}
\end{resolution}

%------------------------------------------------------------------------------%
\begin{resolution}{Norma}
  \begin{align*}
    \onslide<+->{\norm{\vec{v}\times\vec{w}}}
    \onslide<+->{& = \norm{\vetortri{-2}{-5}{11}} \\}
    \onslide<+->{& = \sqrt{(-2)^2 + (-5)^2 + 11^2} \\}
    \onslide<+->{& = \sqrt{4 + 25 + 121} \\}
    \onslide<+->{& = \sqrt{150}}
  \end{align*}
\end{resolution}

%------------------------------------------------------------------------------%
\begin{resolution}{Área}
  \[
    \onslide<+->{A_T}
    \onslide<+->{\;=\; \frac{1}{2} \norm{\vec{v}\times\vec{w}}}
    \onslide<+->{\;=\; \frac{1}{2}\sqrt{150}}
  \]
\end{resolution}

%------------------------------------------------------------------------------%
\endinput

%------------------------------------------------------------------------------%
\begin{question}
  Determine a área do paralelogramo determinado por
  \[
    \vec{v} = \vetor{3}{1}
    \qquad
    \vec{w} = \vetor{1}{4}
  \]
\end{question}

%------------------------------------------------------------------------------%
\begin{resolution}{Vetores}
  Consideramos as versões dos vetores em $\R^3$
  \[
    \vec{v} = \vetortri{3}{1}{0}
    \qquad
    \vec{w} = \vetortri{1}{4}{0}
  \]
\end{resolution}

%------------------------------------------------------------------------------%
\begin{resolution}{Produto Vetorial}
  \begin{align*}
    \onslide<+->{\vec{v}\times\vec{w}}
    \onslide<+->{& = \begin{vmatrix}
      \vec{i} & \vec{j} & \vec{k} \\
      3       &  1      &  0      \\
      1       &  4      &  0
    \end{vmatrix}}
    \\
    \onslide<+->{& = \begin{vmatrix}
         1 &  0 \\
         4 &  0
      \end{vmatrix}
      \vec{i}
      \;-\;
      \begin{vmatrix}
        3 &  0 \\
        1 &  0
      \end{vmatrix}
      \vec{j}
      \;+\;
      \begin{vmatrix}
        3 &  1 \\
        1 &  4
      \end{vmatrix}
      \vec{k}}
    \\
    \onslide<+->{
    & = \big( 3 \times 4 - 1 \times 1 \big) \vec{k}}
    \\
    \onslide<+->{& = 11 \vec{k}}
    \onslide<+->{\;=\; \vetortri{0}{0}{11}}
  \end{align*}
\end{resolution}

%------------------------------------------------------------------------------%
\begin{resolution}{Norma e Área}
  \[
    \onslide<+->{\norm{\vec{v}\times\vec{w}}}
    \onslide<+->{\;=\; \norm{\vetortri{0}{0}{11}}}
    \onslide<+->{\;=\; \sqrt{11^2}}
    \onslide<+->{\;=\; 11}
  \]

  \bigskip
  \[
    \onslide<+->{A}
    \onslide<+->{\;=\; \norm{\vec{v}\times\vec{w}}}
    \onslide<+->{\;=\; 11}
  \]
\end{resolution}

%------------------------------------------------------------------------------%
\endinput

%------------------------------------------------------------------------------%
\begin{question}
  Determine a área do paralelogramo determinado pelos vetores
  \[
    \vec{v} = (2, -1, 3)
    \qquad
    \vec{w} = (1, 4, -2)
  \]
\end{question}

%------------------------------------------------------------------------------%
\begin{resolution}{Produto Vetorial}
  \begin{align*}
    \onslide<+->{\vec{v}\times\vec{w}}
    \onslide<+->{& = \begin{vmatrix}
      \vec{i} & \vec{j} & \vec{k} \\
      2       &  -1     &   3     \\
      1       &   4     &  -2
    \end{vmatrix}}
    \\
    \onslide<+->{& = \begin{vmatrix}
        -1 &  3 \\
         4 & -2
      \end{vmatrix}
      \vec{i}
      \;-\;
      \begin{vmatrix}
        2 &  3 \\
        1 & -2
      \end{vmatrix}
      \vec{j}
      \;+\;
      \begin{vmatrix}
        2 & -1 \\
        1 &  4
      \end{vmatrix}
      \vec{k}}
    \\
    \onslide<+->{
    & = \big( (-1)(-2) - 3 \times 4 \big) \vec{i}
      - \big( 2(-2) - 3 \times 1 \big) \vec{j}
      + \big( 2 \times 4 - (-1)1 \big) \vec{k}}
    \\
    \onslide<+->{& = -10 \vec{i} + 7 \vec{j} + 9 \vec{k}}
    \onslide<+->{\;=\; \vetortri{-10}{7}{9}}
  \end{align*}
\end{resolution}

%------------------------------------------------------------------------------%
\begin{resolution}{Norma e Área}
\begin{columns}
\column{0.5\textwidth}
  \begin{align*}
    \onslide<+->{\norm{\vec{v}\times\vec{w}}}
    \onslide<+->{& = \norm{\vetortri{-10}{7}{9}} \\}
    \onslide<+->{& = \sqrt{(-10)^2 + 7^2 + 9^2} \\}
    \onslide<+->{& = \sqrt{100 + 49 + 81} \\}
    \onslide<+->{& = \sqrt{230}}
  \end{align*}
\column{0.5\textwidth}
  \begin{align*}
    \onslide<+->{A}
    \onslide<+->{\;=\; \norm{\vec{v}\times\vec{w}}}
    \onslide<+->{\;=\; \sqrt{230}}
  \end{align*}
\end{columns}
\end{resolution}

%------------------------------------------------------------------------------%
\endinput

%------------------------------------------------------------------------------%
\begin{question}
  Determine o valor de $k$ para que o paralelogramo determinado pelos vetores
  \[
    \vec{v} = \vetortri{k}{2}{1}
    \qquad
    \vec{w} = \vetortri{1}{-1}{2}
  \]
  tenha área igual a
  \(
    3\sqrt{6}
  \)
\end{question}

%------------------------------------------------------------------------------%
\begin{resolution}{Produto Vetorial}
  \begin{align*}
    \onslide<+->{\vec{v}\times\vec{w}}
    \onslide<+->{& = \begin{vmatrix}
      \vec{i} & \vec{j} & \vec{k} \\
      k       &   2     &   1     \\
      1       &  -1     &   2
    \end{vmatrix}}
    \\
    \onslide<+->{& = \begin{vmatrix}
         2 &  1 \\
        -1 &  2
      \end{vmatrix}
      \vec{i}
      \;-\;
      \begin{vmatrix}
        k &  1 \\
        1 &  2
      \end{vmatrix}
      \vec{j}
      \;+\;
      \begin{vmatrix}
        k &  2 \\
        1 & -1
      \end{vmatrix}
      \vec{k}}
    \\
    \onslide<+->{
    & = \big( 2 \times 2 + 1(-1) \big) \vec{i}
      - \big( 2k - 1 \big) \vec{j}
      + \big( -k - 2 \big) \vec{k}}
    \\
    \onslide<+->{& = 3 \vec{i} + (1-2k) \vec{j} + (-k-2) \vec{k}}
    \onslide<+->{\;=\; \vetortri{3}{1-2k}{-k-2}}
  \end{align*}
\end{resolution}

%------------------------------------------------------------------------------%
\begin{resolution}{Norma}
  \begin{align*}
    \onslide<+->{\norm{\vec{v}\times\vec{w}}}
    \onslide<+->{& = \norm{\vetortri{3}{1-2k}{-k-2}} \\}
    \onslide<+->{& = \sqrt{3^2 + (1-2k)^2 + (-k-2)^2} \\}
    \onslide<+->{& = \sqrt{9 + 1^2 -4k + 4k^2 + k^2 + 4k + 4} \\}
    \onslide<+->{& = \sqrt{9 + 1 + 4 -4k + 4k + 4k^2 + k^2} \\}
    \onslide<+->{& = \sqrt{14 + 5k^2}}
  \end{align*}
\end{resolution}

%------------------------------------------------------------------------------%
\begin{resolution}{Área}
  \begin{align*}
    \onslide<+->{A = \norm{\vec{v}\times\vec{w}} & = 3\sqrt{6}                \\}
    \onslide<+->{\sqrt{14 + 5k^2}                & = 3\sqrt{6}                \\}
    \onslide<+->{14 + 5k^2                       & = \left(3\sqrt{6}\right)^2 \\}
    \onslide<+->{14 + 5k^2                       & = 9 \times 6               \\}
    \onslide<+->{5k^2                            & = 54 - 14 }
    \onslide<+->{                                \;=\; 40                     \\}
    \onslide<+->{k^2                             & = \frac{40}{5} }
    \onslide<+->{                                \;=\; 8                      \\}
    \onslide<+->{k                               & = \pm \sqrt{8}}
    \onslide<+->{                                \;=\; \pm 2 \sqrt{2}}
  \end{align*}
\end{resolution}

%------------------------------------------------------------------------------%
\endinput

%------------------------------------------------------------------------------%
\begin{question}
  Considere quatro pontos no espaço
  \begin{align*}
    A = (1, 0, 2) \\
    B = (4, 1, 3) \\
    C = (2, 4, 1) \\
    D = (5, 5, 2)
  \end{align*}
  Verifique se $ABCD$ é um paralelogramo, se for calcule sua área
\end{question}

%------------------------------------------------------------------------------%
\begin{resolution}{Verificação}
  \onslide<+->{Se $ABCD$ é um paralelogramo}

  \onslide<+->{ele pode ser definido pelos vetores}
  \begin{align*}
    \onslide<+->{\vec{v} & = \overrightarrow{AB} \\}
    \onslide<+->{\vec{w} & = \overrightarrow{AD}}
  \end{align*}
  \onslide<+->{com}
  \[
    \onslide<+->{C \;=\; \overrightarrow{AB} + \overrightarrow{AD}}
    \onslide<+->{  \;=\; \vec{v} + \vec{w}}
  \]
\end{resolution}

%------------------------------------------------------------------------------%
\begin{resolution}{Verificação}
\begin{columns}[t]
\column{0.3\textwidth}
  \begin{align*}
    \onslide<+->{\vec{v}}
    \onslide<+->{& = \overrightarrow{AB} \\}
    \onslide<+->{& = B - A \\}
    \onslide<+->{& = (4, 1, 3) - (1, 0, 2) \\}
    \onslide<+->{& = (3, 1, 1)}
  \end{align*}
\column{0.3\textwidth}
  \begin{align*}
    \onslide<+->{\vec{w}}
    \onslide<+->{& = \overrightarrow{AD} \\}
    \onslide<+->{& = D - A \\}
    \onslide<+->{& = (5, 5, 2) - (1, 0, 2) \\}
    \onslide<+->{& = (4, 5, 0)}
  \end{align*}
\column{0.4\textwidth}
  \begin{align*}
    \onslide<+->{\vec{v} + \vec{w}}
    \onslide<+->{& = (3, 1, 1) + (4, 5, 0) \\}
    \onslide<+->{& = (7, 6, 1)             \\}
    \onslide<+->{& \neq (2, 4, 1)}
    \onslide<+->{\;=\; C}
  \end{align*}
\end{columns}

\bigskip
\onslide<+->{Os pontos não formam um paralelogramo}
\end{resolution}

%------------------------------------------------------------------------------%
\endinput


%%------------------------------------------------------------------------------%
%\section{Lista Mínima}
%
%%------------------------------------------------------------------------------%
%\begin{frame}{Lista Mínima}
%  \begin{enumerate}
%    \item
%  \end{enumerate}
%
%  \vfill
%  Atenção: A prova é baseada no livro, não nas apresentações
%\end{frame}

%------------------------------------------------------------------------------%
\end{document}
%------------------------------------------------------------------------------%
